\chapter{The identity component}
\label{pa-ch-pic0}

\begin{definition}[Picard variety]
  \label{pa-def-pic0}
  \dcref{5.1}
  \uses{pa-thm-picard-representability,pa-prop-picard-group}
  \cite[Lem.~5.1]{kleiman-picard}
  Define
  \[
    \operatorname{Pic}^0_{C/k}:=(\operatorname{Pic}_{C/k})^0,
  \]
  the connected component containing the trivial line bundle.
\end{definition}

The degree is locally constant in a flat family, so the identity component is
contained in the degree-zero component.  The Abel map for sufficiently large
degree shows that these components differ only by translation.

\begin{theorem}[Picard variety is an abelian variety]
  \label{pa-thm-pic0-abelian}
  \dcref{5.4,5.11,5.13,5.14,custom}
  \uses{pa-def-pic0,pa-prop-picard-tangent,pa-def-genus}
  \cite[Thm.~5.4 and Cors.~5.13--5.14]{kleiman-picard}
  The scheme \(\operatorname{Pic}^0_{C/k}\) is a smooth proper commutative
  geometrically connected group scheme, hence an abelian variety, and
  \[
    \dim\operatorname{Pic}^0_{C/k}=g(C).
  \]
\end{theorem}

\begin{proof}
  For \(n>2g(C)-2\), Riemann--Roch gives \(H^1(C,\mathcal L)=0\) for every
  line bundle of degree \(n\).  The proper symmetric power \(C^{(n)}\) maps
  surjectively to \(\operatorname{Pic}^n_{C/k}\); translating by a fixed
  degree-\(n\) class proves that \(\operatorname{Pic}^0_{C/k}\) is proper.
  Its tangent space is \(H^1(C,\mathcal O_C)\), while
  \(H^2(C,\mathcal O_C)=0\) removes the obstruction to lifting line bundles
  across square-zero thickenings.  Thus it is smooth at the identity and,
  by translation, everywhere.  The identity-component theorem gives
  geometric connectedness and the group structure, so it is an abelian
  variety of dimension \(g(C)\).
\end{proof}

\begin{proposition}[Degree-zero characterisation]
  \label{pa-prop-pic0-degree}
  \dcref{5.3,custom}
  \uses{pa-def-pic0}
  A geometric point of \(\operatorname{Pic}_{C/k}\) lies in
  \(\operatorname{Pic}^0_{C/k}\) exactly when its line bundle has degree zero.
\end{proposition}

\begin{proof}
  Degree is constant on connected components, so \(\operatorname{Pic}^0\) is
  contained in the degree-zero locus.  Conversely, after an algebraic
  closure, any degree-zero class can be written as
  \([\mathcal O_C(D-nP_0)]\) for \(n\) sufficiently large by
  Riemann--Roch.  It lies in the component of the identity, and descent from
  the algebraic closure gives the assertion over \(k\).
\end{proof}
